\problema{Trapping Rain Water}{42}{src/06-sliding-window-stack/42-trapping-rain-water.cpp}{H}

Dado un arreglo \texttt{height} que representa la altura de columnas de un
histograma (cada columna tiene ancho \texttt{1}), calcula cuánta agua de
lluvia queda atrapada entre las columnas después de llover.

\paragraph{La idea, con un ejemplo de la vida real.} Imagina una fila de
edificios de distintas alturas pegados uno junto a otro, y que llueve
mucho. El agua se acumula entre dos edificios altos que hacen de ``paredes''
para uno más bajo en medio: el nivel del agua sobre cualquier edificio no
puede subir más que el más bajo de los dos edificios más altos que lo
encierran (por la izquierda y por la derecha), porque el agua se derramaría
por el lado más bajo. Entonces, sobre cada edificio, el agua atrapada es la
diferencia entre ese nivel límite y la altura del edificio mismo (si esa
diferencia es negativa, no se atrapa nada ahí).

\begin{center}
\ttfamily
\begin{tabular}{l}
altura: 0 1 0 2 1 0 1 3 2 1 2 1\\
agua atrapada total: 6\\
\end{tabular}
\end{center}

\paragraph{Cómo lo resuelve el código, paso a paso.} En vez de calcular, para
cada columna, el máximo a su izquierda y el máximo a su derecha por
separado (lo cual tomaría espacio extra), usamos \textbf{two pointers} (dos
punteros) que arrancan en los dos extremos del arreglo y se acercan.
\begin{enumerate}
  \item Mantenemos \texttt{leftMax} (el máximo visto hasta el puntero
        izquierdo) y \texttt{rightMax} (el máximo visto hasta el puntero
        derecho).
  \item En cada paso comparamos \texttt{height[L]} contra
        \texttt{height[R]}.
  \item Si \texttt{height[L]} es menor o igual, sabemos con certeza que el
        edificio más alto que limita el agua sobre \texttt{L} por el lado
        derecho es \emph{al menos} \texttt{height[R]}, que ya es mayor o
        igual que \texttt{height[L]}. Entonces el límite real es
        \texttt{leftMax} (el máximo visto por la izquierda), sin necesidad
        de saber el máximo exacto a la derecha. Sumamos
        \texttt{leftMax - height[L]} al total y avanzamos \texttt{L}.
  \item Si es al revés, hacemos lo simétrico con \texttt{R} y
        \texttt{rightMax}.
  \item Repetimos hasta que los punteros se crucen.
\end{enumerate}

\lstinputlisting{src/06-sliding-window-stack/42-trapping-rain-water.cpp}

\complejidad{$O(n)$}{$O(1)$}

\paragraph{¿Qué significa esa complejidad?} Cada paso del ciclo avanza uno
de los dos punteros, y entre los dos recorren el arreglo una sola vez en
total, así que el tiempo es lineal en el tamaño del arreglo. El espacio es
constante porque solo guardamos un puñado de variables (\texttt{leftMax},
\texttt{rightMax}, los dos punteros y el total), sin arreglos auxiliares.

\paragraph{Consejos para la entrevista (Amazon).} Este problema suele
empezar con una solución de dos arreglos auxiliares
(\texttt{maxIzquierda[]}, \texttt{maxDerecha[]}) que ya es \texttt{O(n)}
en tiempo pero también \texttt{O(n)} en espacio; el reto real es notar que
se puede resolver con espacio constante usando two pointers, y explicar
\emph{por qué} es seguro mover el puntero del lado menor sin conocer el
máximo real del otro lado (la clave está en que ese máximo real, aunque
desconocido, ya es mayor o igual al lado que estamos procesando). Casos
límite que conviene mencionar: arreglo vacío o con menos de tres columnas
(no se puede atrapar nada, porque hace falta al menos un valle entre dos
paredes).
